\begin{table}[t]
\centering
\scriptsize
\setlength{\tabcolsep}{3.5pt}
\caption{Safe haven (T1b): Baur--Lucey regressions}
\label{tab:econ_t1}
\begin{tabular}{llrrrrl}
\toprule
Pair & Period & Hedge $b$ & Slope 10\% & Slope 5\% & Slope 1\% & Verdict \\
\midrule
S\&P 500 -- Gold & Train & 0.00 & $-$0.16 & $-$0.20$^{***}$ & 0.01 & strong safe haven \\
 & Test & 0.09$^{*}$ & 0.17 & 0.12 & 0.15 & weak safe haven \\
S\&P 500 -- Bitcoin & Train & 0.45$^{***}$ & $-$0.13 & 0.29 & 1.95$^{***}$ & not a safe haven \\
 & Test & 1.02$^{***}$ & 1.27$^{***}$ & 1.59$^{***}$ & 2.95$^{***}$ & not a safe haven \\
S\&P 500 -- Bond & Train & $-$0.15$^{***}$ & $-$0.17$^{***}$ & $-$0.16$^{***}$ & $-$0.12$^{**}$ & strong safe haven \\
 & Test & 0.15$^{***}$ & $-$0.10 & 0.02 & 0.08 & weak safe haven \\
\bottomrule
\end{tabular}
\par\smallskip
\parbox{\linewidth}{\footnotesize GARCH(1,1) regression of the daily return of the second asset on the S\&P 500 return, with slope shifts on the days when the S\&P 500 return is below its 10\%, 5\% or 1\% quantile in the period (Baur and Lucey, 2010). $b$: average slope (hedge if $b \le 0$). Slope $q$\%: slope on the days below the $q$\% quantile, $b$ plus the shifts of all nested quantiles down to $q$. Strong safe haven: slope significantly negative at some $q$; not a safe haven: significantly positive at some $q$; weak safe haven otherwise. Train 2015--2021, test 2022--2024; thresholds of the triggers from the training period only. Bond = 10-year zero-coupon Treasury (G\"urkaynak--Sack--Wright). $^{*}$, $^{**}$, $^{***}$: significant at 10\%, 5\%, 1\%.}
\end{table}

\begin{table}[t]
\centering
\scriptsize
\setlength{\tabcolsep}{3.5pt}
\caption{Contagion (T2): Forbes--Rigobon test}
\label{tab:econ_t2}
\begin{tabular}{llrrrrrl}
\toprule
Pair & Period & Days & $\rho$ trigger & $\rho$ other & $\rho^{*}$ trigger & $p$ unadj. & Verdict \\
\midrule
S\&P 500 -- Gold & Train & 86 & 0.27 & $-$0.09 & 0.09$^{*}$ & $<$0.001 & interdependence \\
 & Test & 85 & 0.39 & 0.05 & 0.32$^{***}$ & $<$0.001 & contagion \\
S\&P 500 -- Bitcoin & Train & 99 & 0.43 & 0.10 & 0.15 & $<$0.001 & interdependence \\
 & Test & 114 & 0.58 & 0.36 & 0.41 & 0.003 & interdependence \\
S\&P 500 -- Bond & Train & 46 & $-$0.39 & $-$0.36 & $-$0.15$^{*}$ & 0.593 & no contagion \\
 & Test & 83 & 0.08 & 0.14 & 0.06 & 0.694 & no contagion \\
\bottomrule
\end{tabular}
\par\smallskip
\parbox{\linewidth}{\footnotesize Trigger days: the 50-day window ending on the day has at least two co-crash days (tailDependent, the trigger of rule T2). $\rho$: correlation of daily returns on trigger and other days. $\rho^{*}$: trigger correlation corrected for the higher variance of the S\&P 500 (Forbes and Rigobon, 2002); stars from a one-sided Fisher test of $\rho^{*} > \rho$ other (contagion). $p$ unadj.: same test without correction; significant only before correction = interdependence. Train 2015--2021, test 2022--2024; thresholds of the triggers from the training period only. Bond = 10-year zero-coupon Treasury (G\"urkaynak--Sack--Wright). $^{*}$, $^{**}$, $^{***}$: significant at 10\%, 5\%, 1\%.}
\end{table}

\begin{table}[t]
\centering
\scriptsize
\setlength{\tabcolsep}{3.5pt}
\caption{Decoupling in calm markets (T3b): correlation on calm days}
\label{tab:econ_t3}
\begin{tabular}{llrrrl}
\toprule
Pair & Period & Days & $\rho$ calm & $\rho$ other & Verdict \\
\midrule
S\&P 500 -- Gold & Train & 329 & $-$0.07 & 0.01 & decoupled \\
 & Test & 74 & 0.31$^{***}$ & 0.09 & correlated (+) \\
S\&P 500 -- Bitcoin & Train & 296 & $-$0.08 & 0.19 & decoupled \\
 & Test & 177 & 0.14$^{*}$ & 0.46 & decoupled \\
S\&P 500 -- Bond & Train & 338 & $-$0.20$^{***}$ & $-$0.37 & correlated ($-$) \\
 & Test & 0 & -- & -- & not testable \\
\bottomrule
\end{tabular}
\par\smallskip
\parbox{\linewidth}{\footnotesize Calm days: the 50-day window ending on the day has low volatility for both series (bothCalm, the trigger of rule T3b). Stars test zero correlation on calm days (Fisher); decoupled = not significant. Train 2015--2021, test 2022--2024; thresholds of the triggers from the training period only. Bond = 10-year zero-coupon Treasury (G\"urkaynak--Sack--Wright). $^{*}$, $^{**}$, $^{***}$: significant at 10\%, 5\%, 1\%.}
\end{table}

\begin{table}[t]
\centering
\scriptsize
\setlength{\tabcolsep}{3.5pt}
\caption{Lead under stress (T4): predictive regressions}
\label{tab:econ_t4}
\begin{tabular}{llrrrl}
\toprule
Pair & Period & Lag $k$ & Stress days & $\beta+\gamma$ & Verdict \\
\midrule
S\&P 500 -- Gold & Train & 5 & 136 & 0.036 & no lead \\
 & Test & 5 & 150 & 0.001 & no lead \\
S\&P 500 -- Bitcoin & Train & 1 & 136 & $-$0.604$^{***}$ & leads ($-$) \\
 & Test & 1 & 150 & $-$0.194 & no lead \\
S\&P 500 -- Bond & Train & 1 & 134 & 0.063$^{*}$ & no lead \\
 & Test & 1 & 148 & 0.045 & no lead \\
\bottomrule
\end{tabular}
\par\smallskip
\parbox{\linewidth}{\footnotesize Regression of the daily return of the second asset on its own lag and on the S\&P 500 return $k$ days earlier, with slope shift $\gamma$ on stress days (highVol of the S\&P 500, the trigger of rule T4); HAC standard errors (10 lags). $\beta+\gamma$: lead coefficient under stress. $k$ = lag fitted by the rule (Gold 5, Bitcoin 1; Bond 1). Train 2015--2021, test 2022--2024; thresholds of the triggers from the training period only. Bond = 10-year zero-coupon Treasury (G\"urkaynak--Sack--Wright). $^{*}$, $^{**}$, $^{***}$: significant at 10\%, 5\%, 1\%.}
\end{table}

\begin{table}[t]
\centering
\scriptsize
\setlength{\tabcolsep}{3.5pt}
\caption{Macro-conditioned co-movement (T7): conditional correlation}
\label{tab:econ_t7}
\begin{tabular}{llrrrrl}
\toprule
Pair & Period & Share high & $\rho$ low & $\rho$ high & Difference & Verdict \\
\midrule
S\&P 500 -- Gold & Train & 0.25 & $-$0.10 & 0.01 & 0.12$^{*}$ & no effect \\
 & Test & 0.49 & 0.14 & 0.08 & $-$0.06 & no effect \\
S\&P 500 -- Bitcoin & Train & 0.25 & 0.04 & 0.23 & 0.19$^{***}$ & higher in high infl. \\
 & Test & 0.49 & 0.26 & 0.41 & 0.15$^{**}$ & higher in high infl. \\
S\&P 500 -- Bond & Train & 0.25 & $-$0.35 & $-$0.21 & 0.14 & no effect \\
 & Test & 0.49 & 0.10 & 0.15 & 0.05 & no effect \\
\bottomrule
\end{tabular}
\par\smallskip
\parbox{\linewidth}{\footnotesize Regression of GARCH(1,1)-standardised daily returns of the second asset on those of the S\&P 500, with slope shift on high-inflation days; HAC standard errors (10 lags). High inflation: centred 3-month average of m/m CPI inflation above its 2015--2021 75th percentile (3.61\%), as in the rules (the trigger of rule T7). Share high: share of high-inflation days. $\rho$ low/high: correlation in low/high inflation; difference = $\rho$ high $-$ $\rho$ low. Train 2015--2021, test 2022--2024; thresholds of the triggers from the training period only. Bond = 10-year zero-coupon Treasury (G\"urkaynak--Sack--Wright). $^{*}$, $^{**}$, $^{***}$: significant at 10\%, 5\%, 1\%.}
\end{table}
